\documentclass{comp450}
\usepackage[T1]{fontenc}

\title{Project 3: Barking up a Random Tree}
\author{Darshon Singh (ds161) and Winson Lin (wl85)}

\setlength{\parskip}{0.5\baselineskip}

\begin{document}
\maketitle

\section*{Problem}

We implemented RTP in C++ using OMPL to plan motions for the seven-joint Panda arm in two custom 3D scenes. We compared it with RRT, RRT-Connect and PRM using repeated clutter-scene trials, measuring planning time, path length, smoothness, sampled minimum clearance and exact-solution success rate.

\section*{Panda arm and its joint space}

The Panda has seven revolute joints with a fixed base and gripper. A configuration is $q=(q_1,\dots,q_7)$ in the bounded joint box $\mathcal{C}=\prod_{i=1}^{7}[q_i^{\min},q_i^{\max}]\subset\mathbb{R}^7$. URDF limits are in radians:

\begin{center}
\small
\begin{tabular}{cccccccc}
\toprule
joint & 1 & 2 & 3 & 4 & 5 & 6 & 7 \\
\midrule
min & $-2.97$ & $-1.83$ & $-2.97$ & $-3.14$ & $-2.97$ & $-0.09$ & $-2.97$ \\
max & $2.97$ & $1.83$ & $2.97$ & $0.09$ & $2.97$ & $3.82$ & $2.97$ \\
\bottomrule
\end{tabular}
\end{center}

OMPL uses \texttt{RealVectorStateSpace(7)} and linear joint interpolation. Forward kinematics transforms 59 link-attached collision spheres to check against environment spheres and oriented boxes; self-collision is excluded. The joint-box diagonal is 13.4 rad, giving a maximum sampling interval about 0.067 rad at resolution 0.005.

\section*{RTP implementation}

RTP uses OMPL's state space, validity checker and goal. Its tree is a \texttt{std::vector} of states and parent pointers. Each iteration selects a vertex uniformly, samples a target uniformly or from the goal with probability equal to the goal bias, and validates the full connection with \texttt{checkMotion()}. A valid target is cloned and attached to the selected vertex. This is RTP, with no nearest-neighbor selection or range-limited extension. Parent pointers recover the path; on termination the closest tree vertex gives an approximate solution. \texttt{clear()} and the destructor release stored states and nodes, and \texttt{getPlannerData()} exports the labeled tree.

Fresh OMPL 2.0.2 builds on macOS and Ubuntu Linux (ARM64) passed 25 clear/reuse cycles, exact/approximate and continued solving, trivial/invalid starts, tree labels/edges, and balanced state allocations after destruction. An Apple-only CMake visibility fix restored goal sampling in Python; all four custom binding smoke trials then reached the goal.

\section*{Custom scenes and RTP paths (Exercise 2)}

Wall (\texttt{wall.yaml}): a table with a wall rotated $25^\circ$ around $z$, a hanging ball at 0.78 m, and sphere/box clutter. Both queries move from the left to the right: approximate hand-sphere centroids go from $(.46,.37,.27)$ to $(.62,-.29,.39)$ m (q1) and $(.54,.42,.39)$ to $(.51,-.23,.52)$ m (q2). Direct joint interpolation hits the wall.

Window (\texttt{window.yaml}): a frame at $x=.45$ m, with a roughly $.6\times.5$ m opening, a $15^\circ$ tilted post, and a shelf/ball behind it. Queries start to either side and end through the frame at approximately $(.56,.08,.44)$ and $(.62,.08,.50)$ m. Goals have 4--5 cm clearance, but direct edges hit the frame. All endpoints satisfy joint limits and are collision-free; exact seven-angle configurations are in the YAML files.

\begin{figure}[t]
\centering
\subfigure[wall q1 start/goal]{\includegraphics[width=0.24\linewidth,trim=0 55 0 35,clip]{wall_query1_scene.png}}
\subfigure[wall q1 RTP path]{\includegraphics[width=0.24\linewidth,trim=0 55 0 35,clip]{wall_query1_path.png}}
\subfigure[wall q2 start/goal]{\includegraphics[width=0.24\linewidth,trim=0 55 0 35,clip]{wall_query2_scene.png}}
\subfigure[wall q2 RTP path]{\includegraphics[width=0.24\linewidth,trim=0 55 0 35,clip]{wall_query2_path.png}}
\subfigure[window q1 start/goal]{\includegraphics[width=0.24\linewidth,trim=0 55 0 35,clip]{window_query1_scene.png}}
\subfigure[window q1 RTP path]{\includegraphics[width=0.24\linewidth,trim=0 55 0 35,clip]{window_query1_path.png}}
\subfigure[window q2 start/goal]{\includegraphics[width=0.24\linewidth,trim=0 55 0 35,clip]{window_query2_scene.png}}
\subfigure[window q2 RTP path]{\includegraphics[width=0.24\linewidth,trim=0 55 0 35,clip]{window_query2_path.png}}
\caption{Custom scenes in Viser. Blue is the start, green the goal, faded orange the RTP waypoints, red the gripper path. RTP used goal bias 0.05 and a 10 s limit.}
\label{fig:scenes}
\end{figure}

The four saved exact custom runs use goal bias 0.05 and a 10 s limit. Their times are 0.676, 0.201, 0.117 and 0.603 s (wall q1/q2, window q1/q2), lengths 30.77, 9.61, 11.80 and 24.46 rad, and waypoint counts 8, 4, 4 and 7. Straight-line joint distances are 1.20--1.63 rad, lower bounds rather than feasible paths. RTP makes large swings and gripper loops because its full-length edges can span the joint box. Sampled minimum clearance is 0.49--3.75 cm; denser checks at intervals at most 0.002 rad remain positive, with window q2 decreasing from 1.40 to 1.22 cm. The saved window q2 seed-2 run is approximate after 10.065 s: the start is the closest of 829 vertices. We make no repeated-trial custom-scene success claim because the supporting per-trial records are unavailable. Fresh smoke checks are separate from these original Linux measurements.

\section*{Benchmarking RTP, RRT, RRT-Connect and PRM (Exercise 3)}

\textbf{Setup.} The 20 configurations in \texttt{configs/planners.json} comprise 5 RTP biases, 6 RRT bias/range pairs, 5 RRT-Connect ranges and 4 PRM neighbor counts (Table~\ref{tab:bench}). Each database contains 500 runs: five queries, five trials per setting. The main clutter run uses a common 5 s limit; additional databases use clutter at 1/20 s and shelf at 5 s. Each invocation initializes OMPL with seed 1000, then generators advance across trials; different budgets are not necessarily paired samples. The summary CSVs and plots match the databases. Time includes every run; path quality uses exact solutions only and can favor settings that fail on difficult trials. Clutter's direct distances, 1.75--2.65 rad, are path-length lower bounds.

\begin{table}[p]
\centering
\footnotesize
\setlength{\tabcolsep}{3pt}
\begin{tabular}{lcccc|cccc}
\toprule
 & \multicolumn{4}{c|}{exact solutions (of 25)} & \multicolumn{4}{c}{clutter, 5 s, medians} \\
setting & 1 s & 5 s & 20 s & shelf 5 s & time (s) & length (rad) & smoothness & min clear. (cm) \\
\midrule
RTP gb=0.01 & 10 & 24 & 25 & 21 & 1.119 & 24.43 & 0.73 & 5.4 \\
RTP gb=0.05 & 20 & 25 & 25 & 25 & 0.240 & 15.92 & 0.67 & 5.9 \\
RTP gb=0.1 & 25 & 25 & 25 & 25 & 0.199 & 17.12 & 0.66 & 8.1 \\
RTP gb=0.2 & 25 & 25 & 25 & 25 & 0.076 & 13.30 & 0.43 & 7.7 \\
RTP gb=0.5 & 25 & 25 & 25 & 25 & 0.044 & 10.94 & 0.49 & 6.9 \\
\midrule
RRT gb=0.05 r=0.2 & 17 & 17 & 18 & 14 & 0.104 & 4.24 & 667.67 & 1.2 \\
RRT gb=0.05 r=0.4 & 16 & 16 & 21 & 15 & 0.557 & 4.59 & 79.96 & 1.6 \\
RRT gb=0.01 r=1.0 & 15 & 21 & 23 & 20 & 0.556 & 8.42 & 11.57 & 8.8 \\
RRT gb=0.05 r=1.0 & 17 & 24 & 25 & 18 & 0.215 & 7.36 & 10.69 & 4.1 \\
RRT gb=0.2 r=1.0 & 19 & 22 & 24 & 18 & 0.092 & 6.55 & 10.23 & 2.7 \\
RRT gb=0.05 r=2.5 & 19 & 23 & 25 & 23 & 0.336 & 9.23 & 2.03 & 9.8 \\
\midrule
RRT-Connect r=0.2 & 25 & 25 & 25 & 25 & 0.035 & 3.90 & 549.15 & 0.4 \\
RRT-Connect r=0.4 & 25 & 25 & 25 & 25 & 0.029 & 4.80 & 98.61 & 0.6 \\
RRT-Connect r=1.0 & 25 & 25 & 25 & 25 & 0.019 & 4.44 & 8.87 & 3.6 \\
RRT-Connect r=2.5 & 25 & 25 & 25 & 25 & 0.022 & 5.98 & 1.49 & 4.9 \\
RRT-Connect r=5.0 & 25 & 25 & 25 & 25 & 0.028 & 8.37 & 0.53 & 9.1 \\
\midrule
PRM k=5 & 25 & 25 & 25 & 25 & 0.053 & 7.75 & 0.45 & 8.7 \\
PRM k=10 & 25 & 25 & 25 & 25 & 0.065 & 7.91 & 0.44 & 9.9 \\
PRM k=20 & 25 & 25 & 25 & 25 & 0.084 & 8.72 & 0.44 & 4.6 \\
PRM k=40 & 25 & 25 & 25 & 25 & 0.086 & 8.73 & 0.40 & 3.0 \\
\bottomrule
\end{tabular}
\caption{Exact solutions out of 25 runs (5 clutter queries $\times$ 5 trials) for each time limit and for the shelf scene, and median metrics for the main 5 s clutter run. Time includes all runs, including unsuccessful runs near the 5 s limit. The other three metrics are over exact runs only.}
\label{tab:bench}
\vspace{4pt}
\includegraphics[width=\linewidth]{bench_review_selected_metrics.pdf}
\captionof{figure}{Representative clutter settings at 5 s: RTP $gb=.2$, RRT $gb=.05$, $r=1$, RRT-Connect $r=1$, PRM $k=10$. Time includes all trials; other metrics use exact solutions only. Tick labels show exact counts out of 25. Time and smoothness use log scales.}
\label{fig:metrics}
\end{table}

\textbf{Success rate.} RRT-Connect and PRM reached the goal in all 25 trials for every tested setting and budget, including shelf. At 5 s, RTP had 25/25 exact solutions except at $gb=0.01$ (24/25); RRT ranged from 16/25 to 24/25. At $r=0.2$, RRT improved only from 17/25 to 18/25 when the budget rose from 5 to 20 s, whereas $r=1.0$, $gb=0.05$ rose from 24/25 to 25/25. Longer budgets therefore helped some settings more than others. On shelf, RRT achieved 14/25--23/25 while RTP with $gb\geq0.05$ remained at 25/25. Nearest-vertex extension can repeatedly approach an obstacle from similar directions; RTP's random-vertex connections offer different approaches, and RRT-Connect grows from both endpoints. These are plausible explanations, not diagnoses established by tree-growth measurements.

\textbf{Planning time.} RRT-Connect had median times of 0.019--0.035 s, PRM 0.053--0.086 s, and RTP 0.044--1.119 s. RTP's median fell from 1.119 s at $gb=0.01$ to 0.044 s at 0.5, while its median exact-solution tree size fell from 95.5 to 5 vertices. More frequent goal attempts helped on clutter. RRT at $r=1.0$ improved its median from 0.556 s at $gb=0.01$ to 0.092 s at 0.2, but the latter solved only 22/25 trials. Its mean time was 1.093 s, versus 0.850 s and 24/25 exact for $gb=0.05$. A fast median alone can hide unsuccessful trials.

\textbf{Path length.} RTP's median lengths were 10.94--24.43 rad, compared with 3.90--8.37 for RRT-Connect, 4.24--9.23 for RRT, and 7.75--8.73 for PRM. RTP's unbounded-in-range random edges produce substantial detours. Small extensions gave RRT-Connect its shortest paths at $r=0.2$; larger $r=5.0$ increased the median from 3.90 to 8.37 rad. RRT's shortest-path setting solved only 17/25, so its 4.24 rad median describes an easier surviving subset. PRM stopped after finding a connection through a sparse roadmap; these first solutions do not measure a fully developed roadmap's path quality.

\textbf{Smoothness.} OMPL sums $[2\theta/(a+b)]^2$ over interior waypoints, where $\theta$ is the turn angle and $a,b$ are adjacent edge lengths. Many short segments can therefore score worse than a few long ones: RRT-Connect's median decreased from 549.15 at $r=0.2$ to 0.53 at $r=5.0$. The best medians were PRM $k=40$ (0.40) and RTP $gb=0.2$ (0.43). These are geometric waypoint metrics, not velocity or acceleration measurements; RTP's low score coexists with large arm excursions.

\textbf{Minimum clearance.} Median sampled clearance rose with RRT-Connect range, from 0.36 cm at $r=0.2$ to 9.11 cm at $r=5.0$. RRT ranged from 1.24 to 9.79 cm, RTP from 5.39 to 8.13 cm, and PRM from 2.99 to 9.92 cm. None optimizes clearance. The observed associations with range are empirical and may also reflect which routes were found; they do not guarantee that longer edges are safer. The custom metric samples entire edges, whereas OMPL's separate solution-clearance property averages waypoint clearances. A dense recheck of the saved RRT-Connect q1 example found 0.84 mm penetration despite 0.27 mm positive clearance at the prescribed resolution. Thus an exact result means the goal was reached and the configured discrete validator accepted its edges, not that continuous collision freedom was proved.

\begin{figure}[t]
\centering
\subfigure[RTP gb=0.2, 6.51 rad]{\includegraphics[width=0.24\linewidth,trim=50 260 50 190,clip]{clutter_query1_RTP.png}}
\subfigure[RRT r=1.0, 10.63 rad]{\includegraphics[width=0.24\linewidth,trim=50 260 50 190,clip]{clutter_query1_RRT.png}}
\subfigure[RRT-C r=1, 2.92 rad]{\includegraphics[width=0.24\linewidth,trim=50 260 50 190,clip]{clutter_query1_RRTConnect.png}}
\subfigure[PRM k=10, 7.13 rad]{\includegraphics[width=0.24\linewidth,trim=50 260 50 190,clip]{clutter_query1_PRM.png}}
\caption{Paths for clutter query 1 with our selected settings (5 s limit, seed 1), from \texttt{plan\_arm.py}. Same colors as Figure~\ref{fig:scenes}.}
\label{fig:clutterpaths}
\end{figure}

\newpage
\textbf{Best settings and trade-offs.} Table~\ref{tab:best} identifies each planner's best observed median setting at 5 s on clutter; the full values and success counts are in Table~\ref{tab:bench}. Across planners, RRT-Connect $r=1.0$ was fastest (0.019 s), $r=0.2$ shortest (3.90 rad), PRM $k=40$ smoothest (0.40), and PRM $k=10$ had greatest sampled clearance (9.92 cm). RRT-Connect $r=1.0$ balanced 25/25 exact solutions, 4.44 rad paths and 3.64 cm median clearance, but its representative path shows that an aggregate median does not describe every run. Figure~\ref{fig:clutterpaths} contrasts its short route with RRT's 12-waypoint detour and RTP/PRM's few long edges.

\begin{table}[h]
\centering
\small
\begin{tabular}{lcccc}
\toprule
planner & time & length & smoothness & clearance \\
\midrule
RTP ($gb$) & 0.5 & 0.5 & 0.2 & 0.1 \\
RRT ($gb,r$) & $(.2,1)$ & $(.05,.2)$ & $(.05,2.5)$ & $(.05,2.5)$ \\
RRT-Connect ($r$) & 1 & .2 & 5 & 5 \\
PRM ($k$) & 5 & 5 & 40 & 10 \\
\bottomrule
\end{tabular}
\caption{Best observed setting within each planner, using medians at 5 s on clutter. Time includes failures; path quality is conditional on exact success.}
\label{tab:best}
\end{table}

For RTP, $gb=0.5$ gave the best time and length, but smoothness favored 0.2 and clearance 0.1. Increasing bias therefore did not improve every metric, and the length trend was not monotonic (15.92 rad at 0.05 versus 17.12 at 0.1). We did not sweep goal bias in the window scene, so cannot infer its best setting there. Small ranges traded shorter paths for lower clearance and larger smoothness scores in RRT-Connect. Rankings describe these 25 pooled trials, without confidence intervals or universal optimality. Resolution and raw, unsmoothed paths were fixed. Original native logs and Planner Arena evidence are unavailable; comparisons use the retained databases and locally generated plots.

\section*{Difficulty, time spent and contributions}

Exercise 1 (implementing RTP): difficulty 5/10, about 6 hours. Exercise 2 (scenes and paths): 6/10, about 7 hours. Exercise 3 (benchmarks and analysis): 6/10, about 7 hours. The hardest part was making the window scene hard but still solvable, and figuring out why RRT kept failing in a scene the other planners found easy.

Darshon Singh (Person A) implemented and debugged RTP and led algorithm and performance analysis. Winson Lin (Person B) was responsible for custom environments, visualization, benchmark configuration and experiments, and performed the subsequent results audit, reproducibility checks, documentation, and report revisions. Testing, interpretation, report preparation, and submission preparation were shared.

\end{document}
